\section{Okounkov bodies of b-divisors}\label{sec:Okounkovbdiv}

Let $X$ be a connected compact Kähler manifold of dimension $n$, let $T$ be a closed positive $(1,1)$-current on $X$ with $\vol T>0$.

Fix a smooth flag $Y_{\bullet}$ on $X$.

\begin{theorem}\label{thm:pobbd}
     The partial Okounkov body $\Delta_{Y_{\bullet}}(T)$ admits the following expression:
    \begin{equation}\label{eq:DeltaasHlim}
    \Delta_{Y_{\bullet}}(T)=\nu_{Y_{\bullet}}(T)+\lim_{\pi\colon Z\rightarrow X} \Delta_{Y_{\bullet}}\Bigl( \{\Reg \pi^*T\}\Bigr),
    \end{equation}
    where $\pi$ runs over the directed set of projective bimeromorphic morphisms to $X$ with $Z$ normal.
\end{theorem}
Here the limit is a Hausdorff limit. Recall that $\nu_{Y_{\bullet}}(T)$ is defined in \cref{def:valcurr}. The notation $\Delta_{Y_{\bullet}}\left( \{\Reg \pi^*T\}\right)$ requires an explanation: Take a modification $\pi'\colon Z'\rightarrow X$ dominating $\pi\colon Z\rightarrow X$ through a map $h\colon Z'\rightarrow Z$, so that $Y_{\bullet}$ admits a lifting $(W_{\bullet},g)$ to $Z'$. Such a $Z'$ exists by \cref{thm:liftableflag}. Then we define
\[
\Delta_{Y_{\bullet}}\Bigl( \{\Reg \pi^*T\}\Bigr)\coloneqq \Delta_{W_{\bullet}}\Bigl( h^*\{\Reg \pi^*T\}\Bigr)g^{-1}.
\]
It follows from \itemref{itm:thm:Okounkovtranmain:3} that this definition is independent of the choice of $\pi'$.


This theorem suggests that we define
\begin{equation}\label{eq:Okoubodbdiv}
\Delta_{Y_{\bullet}}\Bigl( \mathbb{D}(T) \Bigr)\coloneqq \lim_{\pi\colon Z\rightarrow X} \Delta_{Y_{\bullet}}\Bigl( \{\Reg \pi^*T\}\Bigr).
\end{equation}
Then one could rewrite \eqref{eq:DeltaasHlim} as
\[
\Delta_{Y_{\bullet}}(T)=\Delta_{Y_{\bullet}}\Bigl( \mathbb{D}(T) \Bigr)+\nu_{Y_{\bullet}}(T),
\]
which formally resembles and extends \eqref{eq:DeltaanaW0}.
\begin{remark}\label{rmk:bdiv-sec-11-3:L30}
One should be able to prove the existence of the limits like \eqref{eq:Okoubodbdiv} over other base fields, at least after assuming the existence of resolution of singularities. If so, one would get an interesting extension of the theory of partial Okounkov bodies.
\end{remark}

\begin{lemma}\label{lma:valuationT} Let $T$ be a closed positive $(1,1)$-current on $X$. Then we have
    \begin{equation}\label{eq:nuTaslimit}
    \lim_{\pi\colon Z\rightarrow X} \nu_{Y_{\bullet}}\Bigl(\Sing_Z(\pi^*T)\Bigr)=\nu_{Y_{\bullet}}(T),
    \end{equation}
    where $\pi$ runs over the directed set of projective bimeromorphic morphisms to $X$ with $Z$ normal.
\end{lemma}
Here $\Sing_Z(\pi^*T)$ denotes the divisorial part of $\pi^*T$ in Siu's decomposition, namely
\[
\Sing_Z(\pi^*T)=\pi^*T-\Reg(\pi^*T)=\sum_{D\subseteq Z} \nu(T,D)[D],
\]
where $D$ runs over all prime divisors on $Z$.
The valuation of currents is understood as in \cref{def:modelvalcurr}.

\begin{proof}
Let us write $\nu=\nu_{Y_{\bullet}}$ for simplicity. For the purpose of the proof,
let us write $\mathcal{D}_X$ for the directed set of projective bimeromorphic
morphisms $\pi\colon Z\rightarrow X$ with $Z$ normal.


We first observe that $\nu(\Sing_Z(\pi^*T))$ is increasing in $\pi\colon Z\rightarrow X$ with respect to the lexicographic order. Furthermore,
\[
    \nu\left(\Sing_Z(\pi^*T)\right)\leq_{\lex} \nu(\pi^*T)=\nu(T),
\]
where the equality follows from \cref{lma:modelvalcurrbir}.

Given $\pi\colon Z\rightarrow X$, let $W_1$ denote the strict transform of $Y_1$ in
$Z$. The restriction
\[
    \pi_1\colon W_1\rightarrow Y_1
\]
is bimeromorphic. Let $\widetilde{W}_1$ be the normalization of $W_1$, and let
\[
    \widetilde{\pi}_1\colon \widetilde{W}_1\rightarrow Y_1
\]
be the induced morphism. Thus we have a commutative diagram
\[
\begin{tikzcd}
\widetilde{W}_1 \arrow[r] \arrow[d, "\widetilde{\pi}_1"] &
W_1 \arrow[r, hook] \arrow[d, "\pi_1"] &
Z \arrow[d, "\pi"] \\
Y_1 \arrow[r, Rightarrow, no head] &
Y_1 \arrow[r, hook] &
X.
\end{tikzcd}
\]

We argue \eqref{eq:nuTaslimit} by induction on $n\geq 0$. The case $n=0$ is trivial. Assume that
$n>0$ and that the assertion is known in dimension $n-1$.

We may assume that $\nu(T,Y_1)=0$ after replacing $T$ by $T-\nu(T,Y_1)[Y_1]$. By definition,
\[
    \nu(T)=\Bigl(0,\mu\left(\Tr_{Y_1}(T)\right)\Bigr),
\]
where $\mu$ denotes the valuation induced by the truncated flag
\[
    Y_1\supseteq Y_2\supseteq \cdots \supseteq Y_n.
\]

Observe that the induced morphisms $\widetilde{\pi}_1\colon \widetilde{W}_1\rightarrow Y_1$ are cofinal in the directed set $\mathcal{D}_{Y_1}$. Indeed, modifications obtained
by compositions of blow-ups with smooth centers on $Y_1$ are cofinal, and they are realized by blowing up $X$ along the same centers.\footnote{It is in this inductive
step that we are forced to introduce singularities, as $W_1$ is not smooth in general.}

Therefore, by the induction hypothesis applied to $\Tr_{Y_1}(T)$,
\[
    \mu\left(\Tr_{Y_1}(T)\right)
    =
    \lim_{\pi\colon Z\rightarrow X}
    \mu\left(
        \Sing_{\widetilde{W}_1}
        \left(\widetilde{\pi}_1^*\Tr_{Y_1}(T)\right)
    \right)=\lim_{\pi\colon Z\rightarrow X}
    \mu\left(
        \Sing_{\widetilde{W}_1}
        \left(\Tr_{W_1}\left(\pi^*T\right)\right)
    \right),
\]
where the second equality follows from \cref{lma:rescommpullback}.

Here $\pi$ runs over $\mathcal{D}_X$. We will prove the needed lower
approximation for the tail coordinates in the following cofinal form. Fix
$\pi\colon Z\rightarrow X$, an effective divisor
\[
    D_0=\sum_{\ell=1}^N b_\ell F_\ell\textup{ on }\widetilde{W}_1,\quad [D_0]\leq \Sing_{\widetilde{W}_1}
    \left(\Tr_{W_1}\left(\pi^*T\right)\right).
\]
Fix also $\epsilon>0$. It suffices to establish the following assertion: after replacing $\pi$ by a higher model in $\mathcal{D}_X$ and replacing $D_0$ by its pull-back, we have
\begin{equation}\label{eq:valuationTcofinalcomparison}
    \mu\Bigl(\left.\Sing_Z(\pi^*T)\right|_{W_1}\Bigr)
    \geq \mu\Bigl( [D_0]\Bigr)-(\epsilon,\dots,\epsilon).
\end{equation}

Note that we can freely replace $\pi$ by a higher model in $\mathcal{D}_X$ and $D_0$ by its pull-back.

We may therefore assume that $Z$ is smooth, that $W_1$ is smooth, and that $D_0$ is a simple normal crossing divisor. Furthermore, we may assume that there are prime divisors $E_1,\ldots,E_N$ on $Z$ not contained in $W_1$ such that $(E_\ell)|_{W_1}=F_\ell$ for each $\ell=1,\ldots,N$. We choose the realization with sufficiently high normal weight along $W_1$. To see what this means, work near a general point of $F_\ell$ and choose coordinates with $W_1=\{s=0\}$ and $E_\ell=\{t=0\}$, so that $F_\ell=\{s=t=0\}$. Put $u_0=s$. Starting from the pair
\[
    W_{1,0}=\{u_0=0\},\qquad E_{\ell,0}=\{t=0\},
\]
we repeatedly blow up their intersection. Assume inductively that, after $r-1$ steps, we work in the chart
\[
    s=t^{r-1}u_{r-1},
\]
where the strict transform of $W_1$ is $\{u_{r-1}=0\}$ and the current ambient divisor realizing $F_\ell$ is $\{t=0\}$. The next center is $\{u_{r-1}=t=0\}$. In the chart of this blow-up meeting the strict transform of $W_1$, we have $u_{r-1}=tu_r$, hence
\[
    s=t^r u_r.
\]
This proves the asserted chart by induction on $r$. Thus, by continuing this elementary construction, which induces an isomorphism on $W_1$ near the generic point of $F_\ell$, and then taking a common model for the finitely many divisors, we can make the integer $r=r_\ell$ as large as needed below. Concretely, if
\[
    g=\sum_{q\geq 0}s^q g_q(t,z_3,\ldots,z_n),
\]
then
\begin{equation}\label{eq:ordEell}
    \ord_{E_\ell}(g)=\min_{q\geq 0}\left(r_\ell q+\ord_{F_\ell}(g_q)\right)
\end{equation}
for the chosen integer $r_\ell$.

We may add a Kähler form to $T$, since this changes neither the divisorial parts nor the valuations in question. Thus assume that $T$ is a Kähler current. Write $T=\theta+\ddc\varphi$, and take a quasi-equisingular approximation $(\varphi_j)_j$ of $\varphi$ in
$\PSH(X,\theta)$ such that
\[
    T_j\coloneqq \theta+\ddc\varphi_j
\]
is a Kähler current with analytic singularities and $\varphi_j\geq\varphi$.

If $D_0=0$, there is nothing to prove. Otherwise, discarding zero coefficients from
$D_0$, choose positive numbers $\delta_\ell<b_\ell$ so small that
\[
    \sum_{\ell=1}^N \delta_\ell\,\mu([F_\ell]) \leq (\epsilon,\ldots,\epsilon).
\]
By \cref{cor:quasi-equichar}, $(\pi^*T_j)_j$ is a quasi-equisingular approximation of $\pi^*T$. Hence \cref{thm:Lelongcont} implies that, for large enough $j$, we have
\begin{equation}\label{eq:tracecoeffTj}
    \nu\left(\Tr_{W_1}(\pi^*T_j),F_\ell\right)> b_\ell-\frac{\delta_\ell}{2}
\end{equation}
for every $\ell=1,\ldots,N$. Fix such a $j$.



Let $\chi$ be a local potential of $\pi^*T_j$ near a general point of $F_\ell$. Since $T_j$ has analytic singularities, write
\[
    \chi=c\log\left(\sum_a |g_a|^2\right)+O(1).
\]
Moreover, $\nu(T,Y_1)=0$ and $\varphi_j\geq\varphi$, hence $\nu(T_j,Y_1)=\nu(\pi^*T_j,W_1)=0$. Thus not all restrictions $g_a|_{W_1}$ vanish identically. Writing
\[
    g_a=\sum_{q\geq 0}s^q g_{a,q},
\]
the trace of $\chi$ to $W_1$ is
\[
    c\log\left(\sum_a |g_{a,0}|^2\right)+O(1).
\]
Hence, putting
\[
    m_\ell\coloneqq \min_a\ord_{F_\ell}(g_{a,0}),
\]
we have
\[
    c m_\ell
    =
    \nu\left(\Tr_{W_1}(\pi^*T_j),F_\ell\right)
    >b_\ell-\frac{\delta_\ell}{2}
\]
by \eqref{eq:tracecoeffTj}. For the fixed finite set of generators $g_a$, choose the integer $r_\ell$ in \eqref{eq:ordEell} so large that $r_\ell>m_\ell$. Since the coefficients $g_{a,q}$ are holomorphic, $\ord_{F_\ell}(g_{a,q})\geq 0$ whenever $q\geq 1$. Therefore the terms with $q\geq 1$ in \eqref{eq:ordEell} have order $>m_\ell$, while the terms with $q=0$ have minimum $m_\ell$. Thus
\[
    \nu(T_j,E_\ell)
    =
    c\min_a\ord_{E_\ell}(g_a)
    =
    c m_\ell
    >b_\ell-\delta_\ell.
\]
Since $\varphi_j\geq\varphi$, we have $\nu(T,E_\ell)\geq\nu(T_j,E_\ell)$, and hence
\[
    \nu(T,E_\ell)>b_\ell-\delta_\ell.
\]
Thus,
\[
    \Sing_Z(\pi^*T)
    \geq
    \sum_{\ell=1}^N (b_\ell-\delta_\ell)[E_\ell].
\]
We get
\[
\begin{aligned}
    \mu\Bigl(\left.\Sing_Z(\pi^*T)\right|_{W_1}\Bigr)
    &\geq
    \sum_{\ell=1}^N
    (b_\ell-\delta_\ell)\mu\left(\Tr_{W_1}[E_\ell]\right)       \\
    &\geq
    \sum_{\ell=1}^N
    (b_\ell-\delta_\ell)\mu([F_\ell])                                \\
    &\geq
    \mu([D_0])-(\epsilon,\ldots,\epsilon).
\end{aligned}
\]
This proves \eqref{eq:valuationTcofinalcomparison}, and hence the lemma.
\end{proof}

\begin{proof}[The proof of \cref{thm:pobbd}]
We shall write $\nu=\nu_{Y_{\bullet}}$.

We argue by induction on $n$. The case $n=0$ is of course trivial. Let us assume that $n>0$ and the result is known in dimension $n-1$.

We may replace $T$ by $T-\nu(T,Y_1)[Y_1]$, so that we may reduce to the case where $\nu(T,Y_1)=0$.

For any projective bimeromorphic morphism $\pi\colon Z\rightarrow X$ with $Z$ normal, it follows from \cref{thm:Deltapartialint} (which also holds for a normal variety, as can be seen after passing to a resolution) that we have
\[
\Delta_{Y_{\bullet}}\Bigl(\{\Reg \pi^*T\}\Bigr)=\overline{ \Bigl\{\nu(S): S\in \{\Reg \pi^*T\}\textup{ is closed, positive} \Bigr\} }.
\]
If $R$ is a closed positive current in $\{\Reg \pi^*T\}$, then $R+\Sing_Z(\pi^*T)$ is a closed positive current in the pull-back class $\pi^*\{T\}$.

By the bimeromorphic invariance of psh functions \cref{cor:PSH}, there is a closed positive current $S\in\{T\}$ such that
    \[
        \pi^*S=R+\Sing_Z(\pi^*T).
    \]
    In particular $\pi^*S\geq \Sing_Z(\pi^*T)$. Moreover, by \cref{lma:modelvalcurrbir} and the linearity of the valuation \cref{prop:nuvaluationlinear}, we have
    \[
        \nu(S)=\nu(R)+\nu\Bigl(\Sing_Z(\pi^*T)\Bigr).
    \]
    Taking the closure over all such $R$ gives
    \[
    \Delta_{Y_{\bullet}}\Bigl(\{\Reg \pi^*T\}\Bigr)+\nu\Bigl(\Sing_Z(\pi^*T)\Bigr)\subseteq \overline{ \left\{\nu(S): S\in \{T\}, \pi^*S \geq \Sing_Z(\pi^*T) \right\} }.
    \]
We observe that the right-hand side is decreasing with respect to $\pi$, which together with \cref{lma:valuationT} implies that the net of convex bodies $\Delta_{Y_{\bullet}}\left(\{\Reg \pi^*T\}\right)$ for various $Z$ is uniformly bounded. Suppose that $\Delta$ is the limit of a subnet. Then we have
\[
\Delta+\nu(T)\subseteq \overline{ \Bigl\{\nu(S): S\in \{T\}, S \preceq_{\mathcal{I}} T \Bigr\} }.
\]
As shown in \cref{thm:Deltapartialint}, the right-hand side is exactly $\Delta_{Y_{\bullet}}(T)$. So
\[
\Delta+\nu(T)\subseteq \Delta_{Y_{\bullet}}(T).
\]
But observe that both sides have the same volume: by \cref{thm:pobmain} and the definition of the volume of $\mathbb{D}(T)$, any Hausdorff limit $\Delta$ has volume $\vol\mathbb{D}(T)/n!$, which equals $\vol T/n!$ by \cref{thm:DT}; translation by $\nu(T)$ does not change volume. So equality holds.

It follows from the Blaschke selection theorem \cref{thm:Blaschke} that the limit in \eqref{eq:DeltaasHlim} exists and \eqref{eq:DeltaasHlim} holds.
\end{proof}
